\documentclass[10pt]{article}

\usepackage[utf8]{inputenc}

\usepackage[T1]{fontenc}

\usepackage{amsmath}

\usepackage{amsfonts}

\usepackage{amssymb}

\usepackage[version=4]{mhchem}

\usepackage{stmaryrd}

\usepackage{bbold}

\usepackage{graphicx}

\usepackage[export]{adjustbox}

\graphicspath{ {./images/} }



\begin{document}

СПЕКТР ДИНАМИЧЕСКОИ СИСТЕМЫ $\left\{T_{t}\right\}$ с фазовым пірост ранством $X$ й йвариа н т н о ї м е p o ii $\mu-$ общее название для различных спектральных инвариантов и спектральных свойств соответствующеї группы (или полугруппы) унитарных (изометрических) операторов сдвига



\[

\left(U_{t} f\right)(x)=f\left(T_{t} x\right)

\]



в гильбертовом пространстве $L_{2}(X, \mu)$. Для динамич. системы в узком смысле слова - измеримого потока $\left\{T_{t}\right\}$ или каскада $\left\{T^{n}\right\}$ - речь идет о спектральных иввариантах одного нормального оператора: во втором случае - унитарного оператора $\left(U_{T} f\right)=f(T x)$, а в первом - производящего самосопряженного оператора $A$ однопараметрич. группы унитарных операторов $\left\{U_{t}\right\}$ (согласно теореме Стоуна, $U_{t}=e^{i t} A$ ).



"Спектральная" терминология в теории динамич. систем несколько отличается от обычной. Спектр $U_{T}$ (или $\boldsymbol{A}$ ) в обычном смысле, т. е. множество тех $\lambda$, при к-рых оператор $U_{T}-\lambda E$ (или $A-\lambda E$ ) ве имеет ограниченного обратного, для всех практически интересных $T$ и $\left\{T_{t}\right\}$ совпадает с окружностью $|\lambda|=1$ или с $\mathbb{R}$ (см. [1], [2]). Поэтому: a) спектр в обычном смысле не содержит информации о свойствах данной динамич. системы, отличающих ее от других; б) у спектра в обычном смысле слова практически никогда нет изолированных точек, так что спектр является непрерывным (в обычном смысле слова), причем это огять-таки не содержит информации о специфич. свойствах данной системы. ІІо этой причине в теории динамич. систем говорят о в е іл р е ры в но м с пा ект ре в случае, когда у $U_{T}$ или $A$ нет собственных функциї, кроме констант; о д и с к р е т н о м с п е к т р е - когда собственные функции образуют полную систему в $L_{2}(X, \mu)$; о с меIII а н н о м с пі е к т р е-в остальных случаях.



Свойства динамич. системы, к-рые определяются ее С. Д. с., наз. с п е к т р а ль н ы ми. Таковы эргодичность (она эквивалентна тому, чтобы у $U_{T}$ - соответственно у $A$ - собственное значение 1 - соответственно 0 - было однократным) и переяешивание. Имеется полная метрич. классификация эргодич. динамич. систем с дискретным спектром; такая система определяется С. Д. с. с точностью до метрич. изоморфизма [3]. Аналогичная теория построена и для более общих групп преобразованиї, нежели $\mathbb{R}$ и $\mathbb{Z}$ (см. [4]). В некоммутативном случае формулировки усложняются, шричем С. д. с. уже не полностью определяет систему. Если спектр не дискретный, то ситуация гораздо сложнее. N. Jum.: [1] I on e s c u T T u c c e a A. A B H.: Ergodic theory, ternat. conference on dynamical systems in math. physics, p., 1976



\includegraphics[max width=\textwidth, center]{2023_07_09_d70c98e8f6186da8d051g-1(1)}

Ф о м и н С. В., Эргодическая теория, М., 1980; [4] М а К-

к и Г. У., в кн.: А у с л е н д е р Л., Г р и н Л., Х а н Ф.,



\begin{center}

\includegraphics[max width=\textwidth]{2023_07_09_d70c98e8f6186da8d051g-1(3)}

\end{center}



СШЕКТР КОЛЬЦА - окольцованное топологич. пространство Spec $A$, точками к-рого являются простые идеалы $\mathfrak{p}$ кольца $A$ с Зариского топологией на нем (к-рая наз. также с п е к т р а л ь в о й т 0 п о л ог и е й). При этом предполагается, что кольцо $A$ коммутативно и с единицей. Элементы кольца $A$ можно рассматривать как функции на пространстве Spec $A$, полагая $a(\mathfrak{p})=a \bmod \mathfrak{p} \in A / \mathfrak{p}$. Пространство Spec $A$ несет пучок локальных колец $\mathcal{O}$ (Spec $A$ ), называемый с т р у к т у рены м п у ч ком. Для точки $\mathfrak{p} \in \operatorname{Spec} A$ слой пучка $G(\operatorname{Spec} A)$ над $\mathfrak{p}-$ это локализация $A_{\mathfrak{p}}$ кольца $A$ относительно



Любому гомоморфизму колец $\varphi: A \rightarrow A^{\prime}$, переводящему единицу в единицу, отвечает непрерывное отображение $\varphi^{*}: \operatorname{Spec} A^{\prime} \rightarrow \operatorname{Spec} A$. Если $N-$ нильрадикал кольца $A$, то есгественное отображение $\operatorname{Spec} A / N \rightarrow$ $\rightarrow \operatorname{Spec} A$ является гомеоморфизмом топологич. пространств. Для невильпотентного элемента $f \in A$ пусть $D(f)=$ $=$ Spec $A \backslash V(f)$, где $V(f)=\{\mathfrak{p} \in \mathrm{Spec} A \mid f \in \mathfrak{p}\}$. Тогда окольцованные пространства $D(f)$ и Spec $A_{(f)}$, где $A_{(f)}-$ локализация $A$ относительно $f$, изоморфны. Множества $D(f)$ наз. г ла в вы м и о тк р ытым и мвожес т в а м и. Они образуют базис топологич. пространства $\operatorname{Spec} A$. Точка $\mathfrak{p} \in \operatorname{Spec} A$ замкнута тогда и только тогда, когда $\mathfrak{p}$ - максимальный идеал кольца $A$. Сопоставляя точке $\mathfrak{p}$ ее замыкание $\bar{p}$ в $\operatorname{Spec} A$, получают взаимно однозначное соответствие между точками пространства Spec $A$ и множеством замкнутых неприводимых подмножеств в $\operatorname{Spec} A$. Пространство Spec $A$ квазикомпактно, но, как правило, не является хаусдорфовым. Р а з м е р в о с т ь ю пространства $\operatorname{Spec} A$ наз. наибольшее $n$, для к-рого существует цепочка различных замкнутых неприводимых множеств $Z_{0} \subset . . \subset$ $\subset Z_{n} \subseteq \operatorname{Spec} A$.



Многие свойства кольца $A$ можно охарактеризовать в терминах топологич. пространства $\operatorname{Spec} A$. Напр., кольцо $A$ нётерово тогда и только тогда, когда Spec $A-$ нётерово пространство; простравство $\operatorname{Spec} A$ веприводимо тогда и только тогда, когда кольцо $A / N$ является областью целостности; размерность Spec $A$ совпадает с размерностью Крулля кольца $A$ и т. д.



Иногда рассматривают м а к с и маль в ы й с п е к т p Specm $A$ - подпространство пространства Spec $A$, состоящее из замкнутых точек. Для градуированного кольца $A$ рассматривают также п р о е к т и вн ы й с п е к т р $\operatorname{Proj} A$. Если $A=\sum_{n=0}^{\infty} A_{n}$, то точки $\operatorname{Proj} A$ - это простые однородные идеалы $\mathfrak{p}$ кольца $A$



\begin{center}

\includegraphics[max width=\textwidth]{2023_07_09_d70c98e8f6186da8d051g-1(2)}

\end{center}



Лит.: [1] Б у р б а к и Н., Коммутативная алгебра, пер. с франц., М., 1971; [2] Ш а фф а р е в и ч И. Р., Основы алгебраической геометрии, М., 1972. Л. В. Куззљин.



СПЕКТР МАТРИЦЫ - совокупность ее собственных значений. См. также Характеристический многочлен матрицы.



СПЕКТР ПРОСТРАНСТВ - представляющий объект для обобщенной теорши когомологий, введен в [1].



\includegraphics[max width=\textwidth, center]{2023_07_09_d70c98e8f6186da8d051g-1}

вательность топологических (как правило, клеточных) пространств $\left\{M_{n}\right\}_{n=-\infty}^{\infty}$ вместе с отображениями $s_{n}: \Sigma M_{n} \rightarrow M_{n+1}$, где $\Sigma^{\prime}-$ надстройка. С. І. образуют категорию; при этом мо р фи 3 м с п е к т р а $\boldsymbol{M}$ в спектр $\boldsymbol{N}$ - это, грубо говоря, «кофинальная часть» нек-рой функции $f: \boldsymbol{M} \rightarrow \boldsymbol{N}$, к-рая задается семейством отображениї $f_{n}: M_{n} \rightarrow N_{n}$ с $t_{n} \circ \Sigma f_{n}=f_{n+1} \circ s_{n} \quad\left(t_{n}\right.$ относится к $\boldsymbol{N}$ так, как $s_{n}$ к $M$ ). Определяются также гомотопные морфизмы, вводится понятие гомотопич. эквивалентности С. П. и строится гомотопич. категория С. П. [2]. Вводятся Постникова системы С. П.



Н а д с т р о й к о й $\Sigma \boldsymbol{M}$ над С. П. $\boldsymbol{M}$ наз. С. П. $\Sigma \boldsymbol{M}=\left\{N_{n}\right\}=\Sigma M_{n}$. Пусть $\left(\Sigma^{-1} \boldsymbol{M}\right)_{n}=M_{n-1}$, тогда $\Sigma$ II $\Sigma^{-1}$ - гомотогически взаимно обратные функторы, так что в категории С. П. в отличие от категории пространств функтор надстройки обратим, и это обстоятельство делает удобной категорию С. П. Вообще, в категории С. п. все рассуждения, связанные со стабилизацией (напр., построение спектральной последовательности Адамса), приобретают естественный вид.



П р и м е р ы С. П. 1) Для любого пространства $X$ определяется С. п. $\boldsymbol{X}=\left\{M_{n}\right\}$, где $M_{n}=$ \% при $n<0$ и $M_{n}=\Sigma^{n} X$ при $n \geqslant 0$, a $s_{n}$ есть естественное отождествление $\Sigma\left(\Sigma^{n} X\right) \rightarrow \Sigma^{n+1}(X)$. Так, для $X=S^{0}$ возникает спектр сфер $\left\{S^{n}\right\}$.



\begin{enumerate}

  \setcounter{enumi}{1}

  \item Спектр пространств Эйленберга - Маклейна $H(\pi)$ (или $\boldsymbol{E} \boldsymbol{M}(\boldsymbol{\pi})$ ), где $\pi-$ абелева группа. Гомотопич. эквивалентность $\omega_{n}: K(\pi, n) \rightarrow \Omega K(\pi, n+1)$, где $K(\pi, n)-$ Эйленберга - Маклейна пространство, а $\Omega X-$ петель пространство над $X$, задает сопряженное ото-

\end{enumerate}



\end{document}